\documentclass[11pt]{article}
\usepackage[a4paper,margin=25mm]{geometry}
\usepackage{amsmath,amssymb,amsthm,hyperref}
\newtheorem{lemma}{Lemma}\newtheorem{theorem}{Theorem}
\newcommand{\E}{\mathbb E}
\title{Jackson kernels down to the inverse-dimension endpoint scale}
\author{Uniform analytic extension; independently reviewed by all three collaborators}
\date{30 September 2026}
\begin{document}\maketitle
Use the Jackson notation from
\path{finite_jackson_endpoint_kernel.tex}, with fixed integer $r>3$,
$\ell\asymp m$, $\nu\asymp m$, and degree $d=r(\nu-1)$ a sufficiently
small fixed multiple of $m$. Constants below may depend on these
fixed ratios and on $r$, but not on the center $a$.
Define
\[
 b_a=\frac{\ell}{\nu}(\sin\theta_a+\nu^{-1})
       \asymp\sqrt{a/m}+m^{-1},\qquad
 Q_a(x)=\frac{J_{\nu,r}(\theta(x)-\theta_a)
                  +J_{\nu,r}(\theta(x)+\theta_a)}
                 {\ell(\sin\theta_a+\nu^{-1})}.
\]
Here $0\le a\le L$ with fixed $L$, and $0\le x\le\ell$.

\begin{lemma}[Uniform real localization]\label{real}
The polynomial $Q_a$ is nonnegative, of degree at most $d$, and has
supremum at most $C/b_a\le Cm$. Its beta-reference integral is
bounded above and below by positive constants uniformly in $a$.
On every fixed bounded $x$-interval,
\[
 Q_a(x)\le\frac C{b_a}
       (1+\sqrt m\,|\sqrt x-\sqrt a|)^{-2r}.             \tag{1}
\]
It is at least $c/b_a$ on $|x-a|\le cb_a$, intersected with $[0,\ell]$.
If $ma\ge A_0$ for a sufficiently large fixed $A_0$, then on
$x\in[a/8,8a]$ the right side of (1) is comparable to
\[
 E_a(x)=b_a^{-1}(1+|x-a|/b_a)^{-2r}.
\]
Outside $[a/4,4a]$, with fixed changes of these numerical ratios if
needed, the global real-angle envelope is at most
\[
 \frac C{b_a}(1+m\max(x,a))^{-r}.                       \tag{2}
\]
\end{lemma}
\begin{proof}
Polynomiality and positivity are unchanged. The identity
$\ell\sin\theta_a=2\sqrt{a(\ell-a)}$ gives the scale.
For bounded $x,a$, the angle difference is comparable to
$|\sqrt x-\sqrt a|/\sqrt\ell$, and the reflected angle is farther
from zero than that difference. The trigonometric bounds in the
companion note prove (1) and the local lower bound, including $a=0$.
For large $x$ up to $\ell$, both periodic angular distances are
bounded below by $c\sqrt{x/\ell}$ when $x\ge4a$ after increasing
that fixed factor. This proves (2) globally.
To bound the integral, use $dx=(\ell/2)\sin\theta\,d\theta$ and
$\sin\theta\le\sin\theta_a+|\theta-\theta_a|$.
The normalized Jackson kernel has bounded angular integral and
first absolute moment $O(\nu^{-1})$; its reflected term satisfies
the same estimate. Division by
$\ell(\sin\theta_a+\nu^{-1})$ therefore gives a uniform upper bound
for its Lebesgue integral. The beta density is at most two. The
near-diagonal lower bound and the positive beta density on $[0,L+1]$
give the reference lower bound.
\end{proof}

\begin{lemma}[Uniform local differential comparison]\label{local}
Suppose $A_0/m\le a\le a_0$, where $A_0$ is sufficiently large
and the fixed $a_0>0$ is sufficiently small. On $x\in[a/8,8a]$,
\[
 \frac{|Q_a^{(k)}(x)|}{k!}
 \le C E_a(x)\left[\left(\frac C{kb_a}\right)^k
                               +\left(\frac Ca\right)^k\right].
                                                                    \tag{3}
\]
For a tensor row of unmarked mean $x\asymp a$, suppose its full
range satisfies $H\le C a$. After discarding the exponentially rare
subsets with $|v-x|>c a$, the marked and unmarked low-order comparisons
through $J=B\log m$ have errors
\[
 C E_a(x)\frac{H^2}{mb_a^2},\qquad
 C E_a(x)\frac{H^2}{m}(1+b_a^{-1}+Hb_a^{-2}),             \tag{4}
\]
respectively. The discarded subset contribution is smaller than any
fixed inverse power of $m$ after integration, and the constants are
uniform in $a$ in the stated range.
\end{lemma}
\begin{proof}
On $|\zeta|\le ca$, the angle change is at most
$C|\zeta|/\sqrt{\ell a}$, so the companion complex Jackson estimate
gives $|Q_a(x+\zeta)|\le C E_a(x)e^{C|\zeta|/b_a}$.
Use the Cauchy radius $\min(ca,kb_a/C)$ to prove (3).
The coefficient estimate $|a_k|\le(CH\sqrt{k/N})^k$ then gives the
same sums as before. The uncapped series has ratio
$CH/(b_a\sqrt m)\le C\sqrt a$, made small by $a_0$.
The capped series costs $O(H^2/(ma^2))$, at most a constant times
$H^2/(mb_a^2)$ because $ma\ge A_0$.
For the marked derivative of order $k-1$ its capped cost is
$O(H^2/(ma))$, bounded by $O(H^2/(mb_a))$.
These estimates hold through $J$; alternatively choose the fixed
degree ratio small enough to sum the capped series through degree $d$.

On these rows all $p_j$ are at least $1/2$. The success-weighted
subset law has bias $O(H^2/m)$ and an all-tilt subgaussian variance
proxy $CH^2/m$. Thus $|v-x|>ca$ has probability $e^{-cm}$, and
$\E e^{C|v-x|/b_a}$ and the corresponding weighted second moment
are bounded as in the companion note, since
$H/(b_a\sqrt m)\le C\sqrt a$.
Taylor's theorem proves the scalar subset comparisons in (4).
On the rare subsets, global Markov--Cauchy bounds for the polynomial,
whose supremum is $O(m)$, cost only $\exp(O((\log m)^2))$ through
order $J$. This is absorbed by $e^{-cm}$, uniformly even when
$a=A_0/m$.
\end{proof}

For the next lemma write $x=(\ell/M)s$ for the unmarked row mean,
so $x\asymp s$. Use the reviewed sharp balance and endpoint CDF
bounds in the forms
\begin{equation}\label{inputs}
 \max_jz_j\le C(x+m^{-1}),\qquad
 m\,w\{x\le t\}\le C(t+m^{-1})\quad(t\ge0).
\end{equation}
These are established in the sharp endpoint-balance note, also for
the deleted-column tensor. The marked column bound follows by
absorbing its own contribution to the full deficit, as before.

\begin{theorem}[Relative off-annulus differential tails]\label{relative}
Let $A_0/m\le a\le a_0$ and $J=B\log m$.
There are fixed constants $0<c_0<1<C_0$ such that the normalized
unmarked differential sum, with absolute values on its summands,
over rows $x\notin[c_0a,C_0a]$ is at most
\[
 C(ma)^{1/2-r}+m^{-A}.                                  \tag{5}
\]
For the direct marked differential functional its corresponding
absolute tail is at most
\[
 C a(ma)^{1/2-r}+m^{-A}.                                \tag{6}
\]
Here $A$ is any fixed exponent; choose the global high-order cutoff
constants accordingly. The degree ratio is chosen sufficiently small
after $r$. The statements include the scalar zeroth-order terms.
\end{theorem}
\begin{proof}
On small rows $x\le c_0a$, one has deterministically
$v\le(M/N)x\le2x$, so $v$ stays a fixed relative distance below $a$.
On large rows $x\ge C_0a$, split into $v\ge x/2$ and $v<x/2$.
The latter event has uniform-subset probability $e^{-cm}$:
\eqref{inputs} gives a population range $O(x)$ because $mx\ge C_0A_0$,
and the uniform-subset Maclaurin--Hoeffding bound applies.
Its contribution through $J$ is negligible by the same crude global
polynomial derivative bound as in the companion note and $A_B\le1$.

On all remaining subsets, $v$ is relatively separated from $a$.
The far-angle Cauchy argument in the companion note is uniform here:
with radius $R=k\sqrt{v(\ell-v)}/d$, the inverse-angle displacement
is at most $Ck/d$. The complex trigonometric estimate is
\emph{valid for arbitrary displacement}; it loses $e^{C_r k}$ and
retains the envelope at the original real angle. In particular no
assumption $k\ll\sqrt{ma}$ or displacement $\ll\theta_a$ is used.
Thus (2) survives in the derivative bound:
\[
 \frac{|Q_a^{(k)}(v)|}{k!}
 \le \frac C{b_a}(1+m\max(v,a))^{-r}
       \left(\frac{C_r d}{k\sqrt{v(\ell-v)}}\right)^k.    \tag{7}
\]
For $v\le\ell/2$, sharp balance gives
$\sigma^2\le CN(x+m^{-1})v$.
Combining it with (7) bounds the order-$k$ centered term by its
real envelope times
\[
 \left(C_r\frac dm\sqrt{\frac{x+m^{-1}}k}\right)^k.
\]
Summing over $k$ is at most $C\exp(C_r(d/m)^2(x+m^{-1}))$.
The derivative-order-$(k-1)$ series is bounded by the same expression
times $C(x+m^{-1})$, as is the marked factor $|z_i-v|$.
Here one can use $v\le2x$ and maximize the elementary series
$(t/\sqrt k)^k$; its sum is at most $C(1+t)^C e^{Ct^2}$,
and the extra polynomial is absorbed by changing $C_r$.
All envelopes may now be evaluated using $\max(x,a)$.

Average $A_B$ and reserve its exponential decay
$\E_BA_B\le e^{-c x}$. Choosing $d/m$ small absorbs the preceding
exponential. Write $d\nu=m\,dw$; the CDF hypothesis is
$\nu([0,t])\le C(t+m^{-1})$. Stieltjes integration by parts gives,
for $r>3$,
\begin{align*}
 \frac1{b_a}\int e^{-cx}(1+m\max(x,a))^{-r}\,d\nu(x)
      &\le C\frac a{b_a}(ma)^{-r}
        \le C(ma)^{1/2-r},\\
 \frac1{b_a}\int e^{-cx}(x+m^{-1})
                    (1+m\max(x,a))^{-r}\,d\nu(x)
      &\le C\frac{a^2}{b_a}(ma)^{-r}
        \le C a(ma)^{1/2-r}.
\end{align*}
For example, split at $a$, use $ma\ge A_0$ on the first piece,
and integrate $x^{-r}$ or $x^{1-r}$ against a CDF bounded by
$C(x+m^{-1})$ on the second. The global CDF bound in
\eqref{inputs} is used throughout this integration.

Finally subsets with $v>\ell/2$ are handled by
$\sigma^2\le Nv(\ell-v)$ and the global Bernstein--Walsh bound.
Their rows have $s>N/2$, so the success weight gives $e^{-cm}$,
which absorbs $\exp(Cd^2/N)$ for sufficiently small degree ratio.
The cases $v=0$ contribute only their already bounded scalar term.
This proves (5)--(6).
\end{proof}

\paragraph{Scope.}
This note supplies uniform variable-center kernel, local comparison,
and relative tail estimates. The variable-width local-mass and
weighted-maximum assembly is supplied separately in
\path{../../independent_directions/natural_scale_endpoint_profile.tex}.
\end{document}
